\section{Gonality of \texorpdfstring{$X_1(N)$}{X1(N)}: proof of Theorem D}\label{sec:gonality-x1}

\begin{proof}[Proof of Theorem D]
\emph{Upper bounds.} Derickx and van Hoeij \cite[Table~1]{DerickxVanHoeij2014} give, for each $N$, a modular unit
on $X_1(N)$ over $\Q$ whose degree is the value in Theorem~D.

\emph{Lower bounds.} Let $u$ be the value in Theorem~D and $d=u-1$. For $N=48$ see \S\ref{sec:double-cover}. For the other $N$ we apply Method~II to $X_1(N)$ at the prime
$p$ of \cref{tab:D}, which does not divide $N$; \S\ref{sec:prelim} explains why the reduction of the plane model
of \cite{SutherlandX1} computes $X_1(N)_{\Fp}$, and we check its genus. Phase~1 is CurveArith's: the places of
degree at most $R=d-\min(n,d-1)$ and the expansions of a basis of holomorphic differentials at them, with
CurveArith's precisions. Phase~2 is the search of \S\ref{sec:phase2} over the family $\mathcal F_b(d)$. In every
case it finds no function of degree at most $d$. By \cref{prop:covers} and \cref{lem:kernel},
$\gon_{\Fp}X_1(N)_{\Fp}>d$, and by \cref{lem:specialization}, $\gon_\Q X_1(N)\ge u$.
\end{proof}

\begin{table}[ht]\centering\small
\caption{Theorem D: Method II for $X_1(N)$. $n=\lceil N_1/(p+1)\rceil$; places up to degree $R=d-\min(n,d-1)$;
 $b$ the baseline; $\#\mathcal F_b(d)$ the number of divisors searched. Time: Phase 1 (its longest process, on a shared machine) + Phase 2.}\label{tab:D}
\begin{tabular}{rrrrrrrrrrr}\toprule
$N$ & $g$ & $p$ & $d$ & $N_1$ & $n$ & $R$ & places & $b$ & $\#\mathcal F_b(d)$ & time\\\midrule
41 & 51 & 2 & 21 & 20 & 7 & 14 & 2669 & 2 & 2643150 & 171 min + 14 s\\
42 & 25 & 5 & 11 & 12 & 2 & 9 & 277397 & 1 & 3860784 & 648 min + 153 s\\
43 & 57 & 2 & 23 & 21 & 7 & 16 & 8515 & 2 & 17156720 & 367 min + 50 s\\
44 & 36 & 3 & 14 & 15 & 4 & 10 & 9518 & 1 & 3133892 & 49 min + 25 s\\
45 & 41 & 2 & 17 & 12 & 4 & 13 & 1608 & 2 & 265608 & 86 min + 1 s\\
46 & 45 & 3 & 18 & 22 & 6 & 12 & 69995 & 1 & 374417725 & 884 min + 950 s\\
47 & 70 & 2 & 28 & 23 & 8 & 20 & 110354 & 2 & 651609964 & 1813 min + 3565 s\\
49 & 69 & 2 & 20 & 21 & 7 & 13 & 1497 & 2 & 2167431 & 207 min + 13 s\\
50 & 48 & 3 & 14 & 20 & 5 & 9 & 3369 & 2 & 41835 & 161 min + 1 s\\
\bottomrule\end{tabular}\end{table}


The second phase also counts the divisors it tests and compares the total with an independent count of
$\#\mathcal F_b(d)$ computed from the numbers of places of each degree; they agree in every case.
As a check of the whole pipeline, the same programs reproduce $\gon_{\F_2}X_1(29)=11$ (degree $10$ is
excluded after testing $8568$ divisors, and a function of degree $11$ is found). For $X_1(47)$, an earlier run
with a different Phase~1 driver (a different distribution of the places among the processes; in both runs the
expansions are computed by CurveArith) and an earlier build of the Phase~2 program tested
exactly the same number $651{,}609{,}964$ of divisors and also found no function of degree at most $28$. It ran on the
same machine and with the same version of Magma, so it checks our drivers rather than the platform; its
record is in \texttt{certificates/D\_x1\_gonality/runs/x1\_47\_independent}.

\subsection{The level \texorpdfstring{$48$}{48}: a double cover}\label{sec:double-cover}
For $N=48$ Method~II over $\F_5$ or $\F_7$ would need all places of degree up to $12$, which is out of reach.
We use instead the quotient of $X=X_1(48)$ by the diamond operator $\langle 23\rangle$, an involution; the
quotient is $Y=X_{\{\pm1,\pm23\}}(48)$, of genus $13$, and $X$ has genus $37$. Everything is computed on $Y$.

Let $k$ be a finite field of odd characteristic, and $\pi\colon X\to Y$ a morphism of degree $2$ of smooth
curves over $k$, the quotient by an involution $\sigma$. For a function $a$ on $Y$ write $\pi^*a=a\circ\pi$;
these are exactly the $\sigma$-invariant functions on $X$. Fix a function $w$ on $X$ with $\sigma w=-w$; then
$w^2=\pi^*\delta$ for a function $\delta$ on $Y$. Write $\operatorname{div}_Y\delta=B-2H$ with $B$ the sum of
the branch points, so that $c:=\deg H=\tfrac12\deg B=g_X-2g_Y+1$. (For a hyperelliptic curve
$w^2=\delta(x)$ over $Y=\PP^1$, all of this is classical: every function is $a(x)+b(x)w$.)

\begin{lemma}\label{lem:dc}
\begin{enumerate}
\item Every function $f$ on $X$ is uniquely $f=\pi^*a+\pi^*b\cdot w$ with $a,b$ functions on $Y$, and $f$ comes
 from $Y$ if and only if $b=0$.
\item Put $t=2a$ and $n=a^2-b^2\delta$. Then $f+\sigma f=\pi^*t$, $f\cdot\sigma f=\pi^*n$, and for every place
 $P$ of $Y$ the total pole order of $f$ above $P$ is $\deg P\cdot\max(0,-v_P(t),-v_P(n))$. In particular
 $\deg_Xf\ge\deg_Yn$ and $\deg_Xf\ge\deg_Yt$.
\item If $\deg_Xf=d$ and $b\ne0$, there are an effective divisor $E$ on $Y$ of degree $d-c$ and $u\ne0$, $h$ in
 $L_Y(E+H)$ with $f=(\pi^*h+w)/\pi^*u$.
\end{enumerate}
\end{lemma}
\begin{proof}
(1) Take $\pi^*a=(f+\sigma f)/2$ and $\pi^*b=(f-\sigma f)/(2w)$, both $\sigma$-invariant.
(2) The first two identities are direct. Measure the orders of the two values of $f$ above $P$ in units of
$v_P$, say $\alpha,\beta$ (if a single ramified point $Q$ lies above $P$, then $\alpha=\beta=v_Q(f)/2$; if a single
unramified point $Q$ of degree $2\deg P$ lies above $P$, then $\alpha=\beta=v_Q(f)$). The
pole order above $P$ is $\deg P(\max(0,-\alpha)+\max(0,-\beta))$, while $v_P(n)=\alpha+\beta$ and
$v_P(t)\ge\min(\alpha,\beta)$ with equality if $\alpha\ne\beta$. If $\alpha,\beta\ge0$ both sides vanish; if
$\alpha<0\le\beta$ then $v_P(t)=\alpha\le v_P(n)$ and both sides are $-\alpha$; if $\alpha,\beta<0$ then
$v_P(n)<v_P(t)$ and both sides are $-\alpha-\beta$.
(3) Since $\operatorname{div}_Xw=R-\pi^*H$ with $R$ the ramification divisor, for every divisor $A$ on $Y$ the
poles of $f$ are bounded by $\pi^*A$ if and only if $a\in L_Y(A)$ and $b\in L_Y(A-H)$. Take $A=\pi_*$ of the pole
divisor of $f$, of degree $d$, and $E=\operatorname{div}b+A-H\ge0$, $u=1/b$, $h=a/b$.
\end{proof}

\begin{proposition}\label{prop:dc}
$X$ has no nonconstant function of degree at most $m$ if $2\gon_k Y>m$ and, for every effective divisor $E$ on $Y$
of degree $m-c$, every $u\ne0$ and every $h$ in $L_Y(E+H)$, the function $(\pi^*h+w)/\pi^*u$ has degree $>m$.
\end{proposition}
\begin{proof}
Functions with $b=0$ are $\pi^*a$, of degree $2\deg_Ya$. For the others use \cref{lem:dc}(3) and enlarge $E$
by $m-d$ points, which only enlarges $L_Y(E+H)$.
\end{proof}

For $X_1(48)$ we take $k=\F_7$ and $m=15$, so $c=12$ and $E$ has degree $3$. The diamond operator $\langle23\rangle$
is defined over $\Q$ and extends to the smooth model of $X_1(48)$ over $\Z_{(7)}$; since $2$ is invertible, the
quotient commutes with base change, so $Y_{\F_7}=X_1(48)_{\F_7}/\langle23\rangle$ is the reduction of $Y$. We use
MD~Sage's canonical model of $Y$ (coordinates $y_0,\dots,y_{12}$, a basis of $S_2(\Gamma_Y)$ with integral
$q$-expansions, $y_0=q+O(q^2)$), checked to have good reduction at $7$ as in \S\ref{sec:prelim}. Let
$\psi\in S_2(\Gamma_1(48))$ have integral $q$-expansion, nonzero modulo $7$, with $\langle23\rangle\psi=-\psi$, and put
$w=\psi/y_0$, a function on $X_1(48)$ with $\sigma w=-w$. Then $\delta=w^2=\psi^2/y_0^2$ is a ratio of two forms of
weight $8$ on $Y$, $\delta=\mathrm{Nm}(y)/\mathrm{Dn}(y)$, from an identity $\psi^2\,\mathrm{Dn}(y)=y_0^2\,\mathrm{Nm}(y)$
with integral coefficients checked to the Sturm bound. The reduction $\mathrm{Dn}(\bar y)$ is not identically zero on
$Y_{\F_7}$ (the certificate divides by it), so $\bar\delta=\mathrm{Nm}(\bar y)/\mathrm{Dn}(\bar y)$ is a well-defined
function on $Y_{\F_7}$, and $\psi^2\equiv y_0^2\bar\delta$. Reducing modulo $7$, $\bar w$ is a function on
$X_1(48)_{\F_7}$ with $\sigma\bar w=-\bar w$ and $\bar w^2=\bar\delta$, so the function field of $X_1(48)_{\F_7}$ is
$\F_7(Y)(\sqrt{\bar\delta})$; as a check, Magma computes its genus as $37$. \texttt{CAGonality} gives
$\gon_{\F_7}Y=8$, so $2\gon Y=16>15$. There are $1328$ effective divisors $E$ of degree $3$ on $Y_{\F_7}$; each
$L_Y(E+H)$ has dimension $3$, giving $57$ choices of $u$ up to scalars and $49$ of $h$ modulo $u$, and
$3{,}709{,}104$ functions in all. The bound $\deg_Xf\ge\deg_Yn$ of \cref{lem:dc}(2) excludes all but $18$ of
them, and for those the exact degree, computed on $Y$, is at least $16$. Hence $\gon_{\F_7}X_1(48)>15$, and
$\gon_\Q X_1(48)=16$ by \cref{lem:specialization}. The divisors $E$ are split into parts by a key that depends only on $E$ (not on the order in which Magma lists
places), so the parts cover every $E$ exactly once; the longest part took $12.5$ hours. The
degree formula of \cref{lem:dc}(2) was checked against Magma's degree computation on $X$ itself for a sample of
the candidates, and the whole procedure was tested on $X_1(24)$ and $X_1(26)$, where it reproduces the known
gonalities.
